% Tetromino tilings with line clears
% Draft for The Mathematical Gazette (Article).
% Every numerical claim in this paper has a row in the project's verification
% ledger with the command that reproduces it; see the repository.
\documentclass[11pt,a4paper]{article}

\usepackage[margin=2.5cm]{geometry}
\usepackage{amsmath,amssymb,amsthm}
\usepackage{booktabs}
\usepackage{graphicx}
\usepackage[hidelinks]{hyperref}

\theoremstyle{plain}
\newtheorem{theorem}{Theorem}
\newtheorem{lemma}[theorem]{Lemma}
\newtheorem{corollary}[theorem]{Corollary}
\theoremstyle{definition}
\newtheorem{definition}[theorem]{Definition}
\theoremstyle{remark}
\newtheorem{remark}[theorem]{Remark}

\newcommand{\board}{B}
\newcommand{\Tet}{\mathrm{T}}

\title{Tetromino tilings with line clears}
\author{Jaehyeok Kim}
\date{\today}

\begin{document}
\maketitle

\begin{abstract}
Walkup proved in 1965 that a $w \times n$ rectangle can be tiled by T-tetrominoes
exactly when $4 \mid w$ and $4 \mid n$, and a line of later work measured how badly
other rectangles fail. We ask the same question in a \emph{dynamic} setting: pieces
fall one at a time under gravity, and any row that becomes full is deleted. A
\emph{perfect clear} is a play that starts from the empty board and returns to it.

Line clears collapse the static obstruction almost completely. Four of the seven
tetromino types --- I, T, J and L --- admit a perfect clear at every width, and O
admits one at every even width; we determine the minimum number of pieces in each
case, for all widths. The two exceptions are S and Z, which admit no perfect clear at
any width, alone or together. The classification extends to every subset of the seven
types.

The contrast with the static theory is sharpest for T, where Walkup's condition
reappears \emph{inside} the dynamic construction: the part of the construction that
uses no line clears is exactly a static tiling, so its width must be a multiple of
four, and the line clears cover precisely the residues that static tiling forbids.
\end{abstract}

\section{Introduction}

A tetromino is a polyomino with four cells; there are seven of them up to
translation, conventionally named I, O, T, S, Z, J and L. Which rectangles can be
tiled by copies of a single one?

For the T-tetromino the answer is a classical theorem of
Walkup~\cite{walkup1965}: a $w \times n$ rectangle admits a tiling by
T-tetrominoes if and only if $4 \mid w$ and $4 \mid n$. Korn and
Pak~\cite{kornpak2004} later analysed the structure of those tilings.
Rectangles that fail Walkup's condition cannot be tiled at all, so
Hochberg~\cite{hochberg2015} asked how close one can get: the \emph{gap number} is
the least number of monominoes that must be left uncovered. He proved lower bounds
of a strikingly rigid kind --- width~5 forces at least one monomino every five
columns, width~7 one every seven columns --- and showed the gap number of T is one
of $5, 6, 7, 9$; Torres and Vallejo~\cite{torresvallejo2024} later brought the cap
down to six.

All of that is static. The tetrominoes are also the pieces of Tetris, and there the
board is not a passive region: pieces fall under gravity, and a row that becomes
completely full is deleted, with everything above it dropping by one. The natural
question is what those deletions buy.

\begin{definition}
A \emph{perfect clear} at width $w$ is a play that begins with an empty board of
width $w$ and, after finitely many placements and the deletions they trigger, ends
with the board empty again.
\end{definition}

Our main result is that line clears dissolve the static obstruction for five of the
seven pieces and for no others.

\begin{center}
\begin{tabular}{lll}
\toprule
piece & static & dynamic (this paper) \\
\midrule
I & $4 \mid w$, $4 \mid n$ needed for a full tiling & every width, $n = w/4$ or $w$ \\
O & only even widths & every even width, $n = w/2$ \\
T & Walkup: $4 \mid w$ and $4 \mid n$ & every width, $n = w$ \\
J, L & & every width, $n = w/2$, $w$ or $2w$ \\
S, Z & conjectured unbounded gap number & \textbf{never, at any width} \\
\bottomrule
\end{tabular}
\end{center}

The T row is where the two theories meet most sharply, and it is worth stating the
mechanism now, because it is the shape of the whole paper. Our construction for T
(Theorem~\ref{thm:T}) splits the board into column-disjoint \emph{units}. All but the
last unit must fill their rows without triggering any clear --- that is, they must
\emph{statically tile} a $x \times 4$ rectangle, so Walkup forces $4 \mid x$. Widths
$5$, $6$ and $7$ therefore cannot be built this way at all. They are exactly the
widths the final unit handles, and the final unit is the one allowed to use line
clears. The dynamic rule covers precisely the residues the static theorem forbids.

Two further remarks place this work. Hoogeboom and Kosters~\cite{hoogeboom2004}
characterised which board configurations are reachable from empty when \emph{all
seven} piece types are available; with the full set the only obstruction is a
divisibility condition on the number of cells, so the question is interesting only
when the supply of pieces is restricted. That is the case we treat. Separately,
Tetris is hard: deciding how much of a given board can be cleared is
NP-complete~\cite{demaine2003}, and this remains true for a single piece
type~\cite{mithardness2026}. Our question is different in kind --- not the complexity
of clearing a given board, but which piece sets can empty an \emph{empty} board at
all --- and it turns out to have a clean answer.

\section{The model}
\label{sec:model}

We fix the rules precisely, because one of our results is that they matter less than
one might expect.

The board of width $w$ is a finite set $\board \subseteq \{0,\dots,w-1\} \times
\mathbb{Z}_{\geq 0}$ of occupied cells; we write $(r,c)$ for the cell in row $r$ and
column $c$, with row $0$ at the bottom. Play starts from $\board = \emptyset$. A move
chooses a piece type from the available set $P$, chooses one of its rotations and a
horizontal position, and places its four cells as a set $X$, subject to:

\begin{itemize}
\item[(L1)] $X \cap \board = \emptyset$ --- the four cells were empty;
\item[(L2)] $X$ shifted down by one row is not a legal position --- the piece cannot
  fall further;
\item[(L3)] after the placement, every full row is deleted simultaneously, and each
  surviving cell descends by the number of deleted rows below it.
\end{itemize}

Condition (L2) says a piece locks only when supported. It says nothing about how the
piece reached its position. Real implementations differ here: hard drop (choose a
column, fall straight down), soft drop with horizontal movement (allowing tucks), and
rotation systems with wall kicks (allowing spins) permit successively more
placements. Writing $\mathrm{HD} \subseteq \mathrm{SD} \subseteq \mathrm{SRS}$ for
the three sets of reachable placements, our results are insensitive to the choice:

\begin{itemize}
\item every construction in this paper was found under hard drop only, so it remains
  valid in the larger models;
\item every impossibility proof uses only (L1), (L2), (L3) and the shape of the
  pieces, never the path by which a piece arrived, so it holds in all three.
\end{itemize}

Finally, if a perfect clear uses $n$ pieces and deletes $R$ rows in total, then
counting cells gives
\begin{equation}
4n = Rw,
\label{eq:cells}
\end{equation}
since every cell that is ever placed is eventually deleted, and each deletion removes
$w$ cells.

\section{Two tools}

The first tool does almost all of the work.

\begin{lemma}[Column counting]
\label{lem:count}
In a perfect clear at width $w$ that deletes $R$ rows in total, every column receives
exactly $R$ cells.
\end{lemma}

\begin{proof}
A cell never changes column: placement fixes its column, and the descent in (L3) is
vertical. A deletion removes one cell from each of the $w$ columns, since the deleted
row is full. The board ends empty, so every cell placed in column $c$ is eventually
deleted, and deletions remove cells from column $c$ one at a time, once per deleted
row. Hence column $c$ receives exactly $R$ cells.
\end{proof}

Lemma~\ref{lem:count} is purely combinatorial: it uses (L3) and nothing else, so it is
independent of the rotation system. Summing over columns recovers~\eqref{eq:cells}.

To use it we record, for each piece type, its \emph{column profiles}: the tuple of
cells contributed to consecutive columns, one tuple per rotation.

\begin{center}
\begin{tabular}{ll}
\toprule
piece & column profiles \\
\midrule
I & $(1,1,1,1)$, $(4)$ \\
O & $(2,2)$ \\
T & $(1,2,1)$, $(1,3)$, $(3,1)$ \\
S, Z & $(1,2,1)$, $(2,2)$ \\
J, L & $(2,1,1)$, $(1,1,2)$, $(1,3)$, $(3,1)$ \\
\bottomrule
\end{tabular}
\end{center}

If $x_c$ denotes the number of pieces placed with a given profile starting at column
$c$, Lemma~\ref{lem:count} becomes a system of linear equations in the $x_c$, one per
column, each with right-hand side $R$. We refer to it as the \emph{counting system}.
Note that S and Z have the same profiles, as do J and L; counting alone can never
distinguish them. The reason is structural.

\begin{lemma}[Mirror]
\label{lem:mirror}
The board and the rules are symmetric under left--right reflection, which exchanges
J with L and S with Z. Consequently J and L behave identically at every width and
every piece count, as do S and Z.
\end{lemma}

\begin{proof}
Reflecting every placement of a legal play gives a legal play: (L1) and (L2) are
preserved because reflection is a bijection on cells respecting rows, and a row is
full if and only if its reflection is. Reflection maps the rotations of J onto those
of L and the rotations of S onto those of Z.
\end{proof}

\section{Single piece types}

Throughout this section $P$ consists of one type.

\subsection{S and Z}

\begin{theorem}
\label{thm:S}
S alone admits no perfect clear, at any width. The same holds for Z alone.
\end{theorem}

\begin{proof}
Let $h_c$ be the number of horizontal S pieces (profile $(1,2,1)$) whose leftmost
column is $c$, and $v_c$ the number of vertical ones (profile $(2,2)$). The counting
system reads
\[
h_c + 2h_{c-1} + h_{c-2} + 2v_c + 2v_{c-1} = R \qquad (0 \le c \le w-1),
\]
with the convention that a variable vanishes when its piece would leave the board.
Column $0$ gives $h_0 + 2v_0 = R$ and column $1$ gives $h_1 + 2h_0 + 2v_1 + 2v_0 =
R$; subtracting, $h_1 + h_0 + 2v_1 = 0$, so $h_0 = h_1 = v_1 = 0$ and $2v_0 = R$. The
same subtraction at columns $c+2, c+3$ propagates this: by induction $h_c = 0$ for
all $c$, $v_c = 0$ for odd $c$, and $v_c = R/2$ for even $c$.

If $w$ is odd, then $w-1$ is even, so $v_{w-1} = R/2$; but a vertical piece starting
at column $w-1$ would occupy column $w$, so $v_{w-1} = 0$ and $R = 0$. There is no
nonempty play.

If $w$ is even the counting system is consistent, and we need the geometry. Every
piece is a vertical S starting at an even column, so the columns split into
independent blocks $\{2k, 2k+1\}$, no piece straddling two blocks. A vertical S
occupies relative rows $\{1,2\}$ in its left column and $\{0,1\}$ in its right one:
only the \emph{right} column reaches relative row $0$. Since every piece starts at an
even column, every cell in an even column sits at relative row at least $1$ above the
piece's base, and the base is at absolute row $\geq 0$; so no placement ever fills
cell $(0,0)$.

Nor can a cell descend into it. A cell reaches row $0$ only if all the rows below its
current position are deleted, which requires row $0$ itself to be deleted, which
requires $(0,0)$ to be occupied. The two statements --- ``$(0,0)$ is never filled''
and ``row $0$ is never deleted'' --- therefore support each other, and a simultaneous
induction on the moves establishes both.

The first piece placed in the block $\{0,1\}$ rests on the floor, since nothing else
can occupy those columns, and so leaves a cell in row $0$. Row $0$ is never deleted,
and a cell leaves the board only when its row is deleted; that cell is permanent, and
the board never empties. The statement for Z follows from Lemma~\ref{lem:mirror}.
\end{proof}

\subsection{O and I}

\begin{theorem}
\label{thm:O}
O alone admits a perfect clear exactly at even widths, with $n = w/2$ pieces and
$R = 2$.
\end{theorem}

\begin{proof}
With $o_c$ the number of O pieces starting at column $c$, the counting system is
$2o_c + 2o_{c-1} = R$. Column $0$ gives $2o_0 = R$, column $1$ gives $2o_1 = 0$, and
inductively $o_c = R/2$ for even $c$ and $0$ for odd $c$. If $w$ is odd then
$o_{w-1} = 0$ because the piece would leave the board and $o_{w-2} = 0$ because
$w-2$ is odd, so column $w-1$ gives $R = 0$. If $w$ is even, place $w/2$ pieces side
by side: rows $0$ and $1$ fill and are deleted. That $R = 2$ cannot be improved is
part of Theorem~\ref{thm:lower}.
\end{proof}

\begin{theorem}
\label{thm:I}
I alone admits a perfect clear at every width $w$, with $n = w/4$ pieces and $R=1$
when $4 \mid w$, and $n = w$ pieces and $R = 4$ otherwise.
\end{theorem}

\begin{proof}
For the construction, place one vertical I in each column: rows $0$ to $3$ fill and
are deleted, using $n = w$ pieces. If $4 \mid w$, lay $w/4$ horizontal I pieces along
row $0$ instead, deleting one row. The minimality of $R$ in each case follows from
Theorem~\ref{thm:lower} below.
\end{proof}

\subsection{T}

This is the case where the static theory reappears, so we give it in full.

\begin{theorem}
\label{thm:T}
T alone admits a perfect clear at every width $w \geq 4$, and the minimum number of
pieces is exactly $n = w$, with $R = 4$.
\end{theorem}

\begin{proof}[Proof of the lower bound]
By~\eqref{eq:cells} it suffices to show $R \geq 4$. Let $h_c, p_c, q_c$ count the
pieces with profile $(1,2,1)$, $(1,3)$ and $(3,1)$ respectively starting at column
$c$. The counting system is
\[
N_j := h_j + 2h_{j-1} + h_{j-2} + p_j + 3p_{j-1} + 3q_j + q_{j-1} = R .
\]
We use only $j = 0, 1, 2$, which exist as soon as $w \geq 3$.

If $R = 1$, then $N_0 = 1$ forces $(h_0,p_0,q_0) = (1,0,0)$ or $(0,1,0)$; the first
gives $N_1 \geq 2h_0 = 2$ and the second $N_1 \geq 3p_0 = 3$, both exceeding $1$.

If $R = 2$, then $3q_0 \leq 2$ forces $q_0 = 0$ and $(h_0,p_0) \in
\{(2,0),(1,1),(0,2)\}$; then $N_1 \geq 2h_0 + 3p_0$ equals $4$, $5$ or $6$, all
exceeding $2$.

If $R = 3$, the solutions of $N_0 = 3$ are $(3,0,0)$, $(2,1,0)$, $(1,2,0)$,
$(0,3,0)$ and $(0,0,1)$. For the first four, $N_1 \geq 2h_0 + 3q_0 \geq 6$. The
remaining case has $q_0 = 1$ and $h_0 = p_0 = 0$, whence $N_1 = h_1 + p_1 + 3q_1 + 1
= 3$, so $q_1 = 0$ and $(h_1,p_1) \in \{(2,0),(1,1),(0,2)\}$; but then $N_2 \geq
2h_1 + 3p_1 \geq 4 > 3$.

So $R \geq 4$ and $n = Rw/4 \geq w$.
\end{proof}

\begin{proof}[Proof of the construction]
Write the four rotations of T as
\[
\mathsf{U} = \{(0,0),(0,1),(0,2),(1,1)\},\quad
\mathsf{D} = \{(0,1),(1,0),(1,1),(1,2)\},
\]
\[
\mathsf{L} = \{(0,1),(1,0),(1,1),(2,1)\},\quad
\mathsf{R} = \{(0,0),(1,0),(1,1),(2,0)\},
\]
with cells given as (row, column) offsets from the landing position, and write
$\mathsf{X}@c$ for a placement of $\mathsf{X}$ whose leftmost column is $c$.

We decompose the board into units occupying disjoint intervals of columns. Because a
piece never leaves its unit's columns, the landing row of each placement, and whether
a row of that unit is full, depend only on that unit's own contents. Two kinds of
unit are needed, with $R = 4$ throughout:

\smallskip
\noindent\emph{The stacking unit} $S$, of width $4$, is
$\mathsf{U}@0,\ \mathsf{L}@2,\ \mathsf{R}@0,\ \mathsf{D}@1$. With clears switched
off it fills rows $0$ to $3$ of its four columns exactly, using $4$ pieces.

\smallskip
\noindent\emph{The final units} $E_m$, of width $m \in \{4,5,6,7\}$ and $m$ pieces:
\[
\begin{array}{ll}
E_4: & \mathsf{U}@0\ \ \mathsf{L}@2\ \ \mathsf{R}@0\ \ \mathsf{D}@1 \\
E_5: & \mathsf{U}@0\ \ \mathsf{D}@2\ \ \mathsf{R}@0\ \ \mathsf{L}@3\ \ \mathsf{D}@1 \\
E_6: & \mathsf{U}@1\ \ \mathsf{L}@4\ \ \mathsf{D}@2\ \ \mathsf{R}@0\ \ \mathsf{D}@0\ \ \mathsf{D}@3 \\
E_7: & \mathsf{U}@0\ \ \mathsf{U}@3\ \ \mathsf{D}@2\ \ \mathsf{L}@5\ \ \mathsf{D}@4\ \ \mathsf{R}@0\ \ \mathsf{D}@1
\end{array}
\]
With clears switched on, each empties its own columns after deleting four rows.

Now write $w = 4a + m$ with $m \in \{4,5,6,7\}$ and $a \geq 0$; this is possible for
every $w \geq 4$, taking $m \equiv w \pmod 4$ in the stated range. Play $a$ copies of
$S$ in the columns $[0,4a)$, then $E_m$ in the columns $[4a,w)$.

While the copies of $S$ are played, the columns of $E_m$ are empty, so no row of the
board is full and no deletion occurs; each copy of $S$ therefore fills rows $0$ to
$3$ of its own columns, as required. When $E_m$ begins, every column outside it is
full in rows $0$ to $3$, so a row of the board is full exactly when $E_m$ fills that
row within its own columns. Hence $E_m$ evolves exactly as it would on a board of
width $m$ on its own, deleting four rows; those deletions remove the rows of the
copies of $S$ as well, and the board is empty. The piece count is $4a + m = w$,
matching the lower bound.
\end{proof}

\begin{remark}[Where Walkup reappears]
\label{rem:walkup}
The requirement on a stacking unit --- fill rows $0$ to $3$ of $x$ columns with no
deletion --- says exactly that the $x \times 4$ rectangle admits a static tiling by
T-tetrominoes. By Walkup's theorem this forces $4 \mid x$. So widths $5$, $6$ and $7$
admit no stacking unit at all, and those are precisely the widths carried by the
final unit, the one permitted to use line clears. The dynamic rule covers exactly the
residues the static theorem forbids.

Walkup's theorem does not, however, settle whether the $4 \times 4$ tiling can be
\emph{reached} under gravity: support is not a concern of static tiling, and a tiling
in which some piece must be inserted beneath another would be unusable here. That it
can be reached --- by the four placements listed above --- is the one genuinely
dynamic ingredient of the construction.
\end{remark}

\subsection{J and L}

\begin{theorem}
\label{thm:JL}
J alone admits a perfect clear at every width $w \geq 4$, with minimum piece count
\[
n = \begin{cases}
w/2 & (w \equiv 0 \bmod 4), \quad R = 2,\\
w & (w \equiv 2 \bmod 4), \quad R = 4,\\
2w & (w \text{ odd}), \quad R = 8 .
\end{cases}
\]
The same holds for L.
\end{theorem}

\begin{proof}[Sketch]
The construction is again a unit decomposition, with
$\mathsf{A} = (2,1,1)$, $\mathsf{B} = (1,1,2)$, $\mathsf{P} = (1,3)$ and
$\mathsf{Q} = (3,1)$ denoting the rotations of J by their profiles. The units are
\[
\begin{array}{llll}
\mathsf{AB} & \text{width } 4, & R = 2, & \mathsf{A}@0,\ \mathsf{B}@1 \\
\mathsf{PQ} & \text{width } 2, & R = 4, & \mathsf{P}@0,\ \mathsf{Q}@0 \\
\mathsf{PQPQ} & \text{width } 2, & R = 8, & (\mathsf{P}@0,\ \mathsf{Q}@0) \text{ twice} \\
\mathsf{U}_5 & \text{width } 5, & R = 8, & \mathsf{A}@0\ \mathsf{P}@1\ \mathsf{P}@0\ \mathsf{P}@3\ \mathsf{Q}@0\ \mathsf{Q}@3\ \mathsf{A}@2\ \mathsf{A}@0\ \mathsf{P}@3\ \mathsf{B}@1
\end{array}
\]
and the board is built as $\mathsf{AB}^{w/4}$ when $4 \mid w$, as
$\mathsf{PQ}^{w/2}$ when $w \equiv 2 \pmod 4$, and as $\mathsf{U}_5$ followed by
$(w-5)/2$ copies of $\mathsf{PQPQ}$ when $w$ is odd. The verification is the same two
steps as in Theorem~\ref{thm:T} and is a finite check on four units, independent of
$w$. Minimality follows from Theorem~\ref{thm:lower}. L follows by
Lemma~\ref{lem:mirror}.
\end{proof}

\section{Lower bounds at every width}
\label{sec:lower}

The minimality claims above assert, for each piece and each width, that the counting
system has no solution for smaller $R$. For a fixed width this is a finite check. The
following observation upgrades it to all widths at once.

\begin{theorem}
\label{thm:lower}
For each piece type and each $R$, the set of widths at which the counting system is
solvable is eventually periodic in $w$, and the period can be computed. The resulting
minimum values of $R$ are
\begin{center}
\begin{tabular}{ll}
\toprule
piece & minimum $R$ \\
\midrule
I & $1$ if $4 \mid w$, else $4$ \\
O & $2$ if $w$ even; no solution if $w$ odd \\
T & $4$ at every width \\
S, Z & $2$ if $w$ even; no solution if $w$ odd \\
J, L & $2$ if $4 \mid w$; $4$ if $w \equiv 2 \bmod 4$; $8$ if $w$ odd \\
\bottomrule
\end{tabular}
\end{center}
\end{theorem}

\begin{proof}
Sweep the columns from left to right, maintaining as state the vector of
contributions already committed to the columns not yet processed. A profile spans at
most four columns, so the state has at most three components, each an integer in
$[0,R]$ --- any larger value can never be reduced and is dead. The state space is
therefore finite. Processing one column is a fixed relation on states, so it induces
a deterministic map on \emph{sets} of states; iterating it from the initial set
$\{(0,\dots,0)\}$ gives a sequence of subsets of a finite set, which must eventually
cycle.

The rightmost few columns are special --- a profile may no longer fit --- but the
sequence of restricted maps applied there depends only on how many columns remain,
not on $w$. Hence solvability at width $w$ is obtained by applying a fixed tail map
to the $(w - k)$-th term of an eventually periodic sequence, where $k \leq 3$ is the
number of special columns, and is therefore eventually periodic in $w$. Computing the
pre-period and period gives the table. \qedhere
\end{proof}

The entries for S and Z deserve emphasis, because they explain the shape of
Theorem~\ref{thm:S}.

\begin{corollary}
\label{cor:geometry}
At every even width, S alone passes the counting test with $R = 2$. Consequently no
argument based on Lemma~\ref{lem:count} alone can rule out S, and the geometric step
in the proof of Theorem~\ref{thm:S} is unavoidable.
\end{corollary}

Among all seven pieces, this is the only place where counting and geometry part
company: for I, O, T, J and L the counting bound is achieved by an explicit
construction, and for S and Z at odd widths counting already gives a contradiction.
The single gap is S and Z at even widths, and it is exactly the gap
Theorem~\ref{thm:S} fills.

\section{Sets of pieces}

Suppose now that a set $P$ of types is available, and the player may use any of them
in any order.

\begin{theorem}
\label{thm:SZ}
$\{S,Z\}$ together admits no perfect clear, at any width.
\end{theorem}

\begin{proof}[Sketch]
S and Z have the same column profiles, so the counting argument of
Theorem~\ref{thm:S} applies verbatim to their total counts: no horizontal piece is
used, every piece is vertical and starts at an even column, and odd widths are
excluded outright. At even widths the columns again split into independent blocks
$\{2k,2k+1\}$, and the width-$w$ problem reduces to a width-independent question
about one block.

Let $F_0, F_1$ be the sets of filled rows in the two columns of a block and put
$d_j = |F_0 \cap [0,j]| - |F_1 \cap [0,j]|$. Placing a vertical Z at height $r$
increments $d_r$ and $d_{r+1}$; placing a vertical S decrements them; deleting row
$m$ removes $d_m$ from the sequence and changes no other value. One then checks four
invariants of the reachable states: $|F_0| = |F_1|$; the occupied rows form an
initial segment; $|d_j| \leq 1$ for all $j$; and $d_0 \neq 0$ whenever the block is
nonempty. The last one says row $0$ of the block is never full, so it is never
deleted, and the cell the first piece leaves in row $0$ is permanent.
\end{proof}

\begin{theorem}
\label{thm:sets}
A set $P$ of tetromino types admits a perfect clear at width $w$ if and only if
$P$ meets $\{I, T, J, L\}$, or $w$ is even and $O \in P$.
\end{theorem}

\begin{proof}
If $P$ contains one of I, T, J, L, use that type alone
(Theorems~\ref{thm:I},~\ref{thm:T},~\ref{thm:JL}). Otherwise $P \subseteq
\{O,S,Z\}$. All three of these have profile set $\{(1,2,1),(2,2)\}$, so the counting
derivation in the proof of Theorem~\ref{thm:S} applies to their total counts without
change: horizontal pieces are unused, every piece starts at an even column, and at
odd widths $R = 0$. At even widths, if $O \in P$ then $w/2$ copies of O clear two
rows; if not, then $P \subseteq \{S,Z\}$ and Theorems~\ref{thm:S}
and~\ref{thm:SZ} apply.
\end{proof}

\begin{remark}
It is tempting to argue that since $\{S,Z\}$ is the only unusable pair, every set of
three or more types must work. That is false: $\{O,S,Z\}$ has three types and fails
at every odd width. The reason is that O carries a width restriction of its own, so
``only one bad pair'' says nothing about sets built from the three restricted types.
\end{remark}

\section{Restricted supply: the seven-bag randomiser}
\label{sec:sevenbag}

Modern Tetris does not let the player choose the next piece. Pieces arrive in
\emph{bags}: each block of seven consecutive pieces is a permutation of the seven
types. A single slot of \emph{hold} lets the player set one arriving piece aside.

This changes the question, because every construction above repeats one type.

\begin{theorem}
\label{thm:bag}
Under the seven-bag constraint, with the order inside each bag at the player's
choice, the minimum number of pieces in a perfect clear is
\begin{center}
\begin{tabular}{lrrrrrrrrrrr}
\toprule
$w$ & 4 & 5 & 6 & 7 & 8 & 9 & 10 & 11 & 12 & 13 & 14 \\
\midrule
unrestricted & 1 & 5 & 3 & 7 & 2 & 9 & 5 & 11 & 3 & 13 & 7 \\
seven-bag & 1 & 5 & 3 & 7 & 4 & 9 & 10 & 11 & 9 & 13 & 14 \\
\bottomrule
\end{tabular}
\end{center}
\end{theorem}

The pattern is that the constraint is free at odd widths and costly at even ones. At
an odd width the unrestricted minimum already has $R = 4$, and four rows give room to
place seven different types; at an even width the unrestricted optimum is one or two
rows, and a thin region cannot absorb variety. A single row admits only horizontal I
pieces, so $w/4 \geq 2$ is immediately fatal; two rows survive to $w = 8$ and fail
from $w = 10$.

\smallskip
The harder question is what the \emph{order} inside the bags costs, and here we can
report less. Fix $w = 10$ and $n = 10$, the setting of the standard opener. Ten
placements consume at most eleven arrivals, so only the first four pieces of the
second bag matter, giving $P(7,4) = 840$ prefixes to consider for each of the
$7! = 5040$ orders of the first bag.

We sampled $100$ first bags uniformly from the $5040$ with a fixed seed and settled
all $840$ prefixes for each: $84{,}000$ cases. Nine fail. They fall in three first
bags, three apiece:
\[
\begin{array}{ll}
\mathrm{LJTIOZS}: & \mathrm{OLSZ},\ \mathrm{LOSZ},\ \mathrm{LSOZ} \\
\mathrm{LOTJISZ}: & \mathrm{OJZS},\ \mathrm{JOZS},\ \mathrm{JZOS} \\
\mathrm{TLJOISZ}: & \mathrm{OJZS},\ \mathrm{JOZS},\ \mathrm{JZOS}
\end{array}
\]
The other $97$ first bags admit a perfect clear for every prefix, so the failure rate
over the sample is $9/84{,}000 \approx 1.1 \times 10^{-4}$. An earlier first bag,
$\mathrm{TJLISZO}$, was settled separately and fails on four prefixes
($\mathrm{OJZS}$, $\mathrm{OLZS}$, $\mathrm{LOZS}$, $\mathrm{JOZS}$).

A caution about the computation is worth recording, because it changed the answer.
Our first pass used a state cap of $8 \times 10^6$ and reported \emph{no} failures at
all: every one of the nine was sitting in a case that hit the cap, undecided. They
need between $3.8 \times 10^7$ and $5.2 \times 10^7$ states to settle. A capped
search is not a negative result.

Across all thirteen known failures the second-bag prefix is a permutation of
$\{O,S,Z\}$ together with one of $J$ or $L$ --- no exception so far. A tempting
sharper rule, that the failing prefixes end in $\mathrm{ZS}$, is false: the endings
observed are $\mathrm{ZS}$ eight times, $\mathrm{SZ}$ twice, $\mathrm{OS}$ twice and
$\mathrm{OZ}$ once, and the prefix $\mathrm{LSOZ}$ fails after $\mathrm{LJTIOZS}$ but
succeeds after $\mathrm{TJLISZO}$. Failure is a joint property of the two bags, not
of the second one alone.

What the four failing first bags do share is that $S$ and $Z$ both arrive late ---
in position five or later. Fifteen of our hundred samples have that property and
three of them fail; the other eighty-five do not fail at all. The condition also
survived a test made before the computation finished: the last four undecided cases
all lay in first bags where $S$ or $Z$ arrives early, so the condition predicted they
would succeed, and all four did. This is consistent with the mechanism one would
guess --- a vertical I is the main way to resolve an S/Z staircase inside four rows,
and an I arriving after the staircase has been built cannot undo it, which is the
direction Burgiel's theorem~\cite{burgiel1997} points in.

We claim no novelty for this section. The community's own perfect-clear solvers work
in a more general model --- ours is hard-drop only, so an impossibility here is an
impossibility \emph{in our model}, and spins or tucks might rescue some of these
cases. The point is the shape of the answer: perfect clears are almost always
available, failures are rare and structured, and the well-known openers exist to save
search and memorisation rather than to reach positions that are otherwise
unreachable.

\appendix

\section{A parity identity}
\label{app:parity}

The following is folklore among competitive players, where it is used to decide
whether a perfect clear is still reachable; see the Hard Drop wiki~\cite{harddropparity}.
We record it because it has a short proof, generalises to every width, and explains
a natural false conjecture.

Colour the board like a chessboard and set
\[
\varphi(\board) = \sum_{(r,c) \in \board} (-1)^{r+c}.
\]
Every tetromino except T covers two cells of each colour, so placing it leaves
$\varphi$ unchanged; a T covers three of one colour and one of the other, changing
$\varphi$ by $\pm 2$, hence by $2 \bmod 4$.

Consider a deletion event removing full rows $j_1 < \dots < j_k$. For a surviving
cell at row $r$, let $s(r)$ be the number of deleted rows below it; the cell descends
by $s(r)$, multiplying its contribution by $(-1)^{s(r)}$. Put
\[
D = \#\{\text{surviving cells with } s(r) \text{ odd}\},
\qquad
A = \sum_{i=1}^{k} \varphi(\text{row } j_i),
\]
noting that $\varphi$ of a full row is $0$ when $w$ is even and $(-1)^{j}$ when $w$
is odd. Then
\[
\Delta\varphi = -2\!\!\sum_{s(r) \text{ odd}}\!\!(-1)^{r+c} \; - \; A
\equiv -2D - A \pmod 4,
\]
since modulo $4$ the first sum depends only on the parity of the number of cells
involved. Summing over a whole perfect clear, where $\varphi$ starts and ends at $0$,
gives $2\,\#\mathrm{T} \equiv 2\sum D + \sum A \pmod 4$, that is
\[
\#\mathrm{T} \;\equiv\; \sum_{\text{events}} D \;+\; \tfrac{1}{2}\sum_{\text{events}} A
\pmod 2 .
\]
For even $w$ the correction $A$ vanishes and the identity reads: \emph{the parity of
the number of T pieces equals the parity of the total number of cells left above the
deleted rows}. In particular a perfect clear with no deletions at all uses an even
number of T pieces --- which is the false conjecture one arrives at by forgetting
that rows disappear.

\begin{remark}
$D$ must be counted per deleted row, or equivalently as the number of cells with an
odd number of deleted rows below them. Reading it as ``the cells above the deleted
rows'' is correct only for single-row deletions.
\end{remark}

\section*{Computations}

All constructions and all counting claims in this paper are reproducible from the
accompanying repository, which contains for each numerical statement the command that
produces it. The identity of Appendix~\ref{app:parity} was checked against all
$4{,}758{,}805$ perfect clears at widths $4$, $5$, $6$, $7$ and $8$. Of these, the
$260{,}423$ at widths $4$, $5$, $6$ and $8$ were checked by two independent
implementations agreeing on the number of plays in every case; the $4{,}498{,}382$ at
width $7$ have so far been checked by one.

\begin{thebibliography}{99}
\bibitem{walkup1965} D.~W. Walkup, Covering a rectangle with T-tetrominoes,
  \emph{Amer. Math. Monthly} \textbf{72} (1965) 986--988.
\bibitem{hochberg2015} R.~Hochberg, The gap number of the T-tetromino,
  \emph{Discrete Math.} \textbf{338} (2015) 130--138.
\bibitem{kornpak2004} M.~Korn and I.~Pak, Tilings of rectangles with T-tetrominoes,
  \emph{Theoret. Comput. Sci.} \textbf{319} (2004) 3--27.
\bibitem{torresvallejo2024} Torres and Vallejo, Fragments that reduce the gap cap of
  the T-tetromino to six, \emph{Sci. Eng. J.} \textbf{17} (2024) 128--133.
\bibitem{reid2005} M.~Reid, Klarner systems and tiling boxes with polyominoes,
  \emph{J. Combin. Theory Ser. A} \textbf{111} (2005) 89--105.
\bibitem{burgiel1997} H.~Burgiel, How to lose at Tetris, \emph{Math. Gaz.}
  \textbf{81} (1997) 194--200.
\bibitem{brzustowski1992} J.~Brzustowski, \emph{Can you win at Tetris?},
  MSc thesis, University of British Columbia, 1992.
\bibitem{hoogeboom2004} H.~J. Hoogeboom and W.~A. Kosters, How to construct Tetris
  configurations, \emph{Int. J. Intell. Games Simul.} \textbf{3} (2004) 97--105.
\bibitem{demaine2003} E.~D. Demaine, S.~Hohenberger and D.~Liben-Nowell, Tetris is
  hard, even to approximate, in \emph{COCOON 2003}, 351--363.
\bibitem{mithardness2026} MIT Hardness Group, Tetris is hard with just one piece
  type, in \emph{FUN 2026}, LIPIcs \textbf{366}, 32:1--32:16.
\bibitem{baccherini2008} D.~Baccherini and D.~Merlini, Combinatorial analysis of
  Tetris-like games, \emph{Discrete Math.} \textbf{308} (2008) 4165--4176.
\bibitem{harddropparity} Parity, Hard Drop Tetris Wiki.
\end{thebibliography}

\end{document}
